\documentclass[12pt]{article}
\usepackage[margin=1in, paperwidth=8.5in, paperheight=11in]{geometry}
\usepackage{ifpdf,amsmath, amssymb, comment, color, graphicx, stmaryrd,setspace,enumerate,tikz, fancyhdr, wrapfig, textcomp, mathptmx, siunitx}
\usepackage{pdfpages}
\usepackage{enumerate}

\usepackage{hyperref}
\hypersetup{
    colorlinks=true,
    urlcolor=blue,
}

\setlength{\headheight}{14.5pt}
\newcommand{\Q}{\mathbb{Q}}
\newcommand{\R}{\mathbb{R}}
\newcommand{\Z}{\mathbb{Z}}

\newenvironment{red}{\color{red}}{\ignorespacesafterend}


\newcommand{\blue}[1]{\textcolor{blue}{#1}}
\newcommand{\green}[1]{\textcolor{green}{#1}}
\renewcommand{\section}[1]{\begin{center} \textbf{#1} \\\end{center}}
%
\setlength{\parindent}{0in}
%\oddsidemargin=-.25in
\allowdisplaybreaks
\pagestyle{fancy}
\renewcommand{\headrulewidth}{0pt}
\lhead{MATH 202}
\rhead{Fall 2024}
%\lfoot{\copyright\ CLEAR Calculus 2010}
\cfoot{}
\renewcommand{\thefootnote}{*} 
\hyphenpenalty=10000
\begin{document}
%


%\onehalfspacing
\allowdisplaybreaks
%##################################################################
\section{Physics applications of definite integrals -- When is something a $\Delta$?}
Somebody asked me a good question in an exit quiz recently that I was writing a loonnggg explanation in the Canvas comments about, but then I was like, this sucks, let me write it in a document instead, lol.

The question was: When should I label some quantity with a $\Delta$ and when should I not? This is a good question that I don't think I gave a very persuasive answer to in class.

My short answer is, you should use a $\Delta$ when something is \textit{small}, but not when something is \textit{not small}.

For example, in the leaky bucket problem, we slice the journey of the bucket into small pieces. 
\begin{itemize}
\item The work on each slice is \textit{small}, so I call it \(\Delta W\). 
\item The distance across each slice is \textit{small}, so I call it \(\Delta h\).
\item However, the force (aka the weight) is \textit{not small}; it's some reasonable number. It does depend on \(h\), so I'm going to call it \(F(h)\) (not \(\Delta F\)!).
\end{itemize} Therefore \(\Delta W = F(h) \cdot \Delta h\).

To further illustrate, here's the rest. Refer back to the original handout for the situations, and note that you may have used different letters than I'm using here.

\hrulefill

\begin{enumerate}
\item When we slice the second hand of the clock:
\begin{itemize}
\item The kinetic energy of the slice is \textit{small}, \(\Delta K\).
\item The mass of the slice is \textit{small}, \(\Delta m\).
\item The velocity of the slice is \textit{not small} but depends on \(\ell\), \(v(\ell)\).
\end{itemize}
Therefore, \(\Delta K = \frac12\cdot\Delta m\cdot v(\ell)^2\).

\item When we slice the oil slick into rings:
\begin{itemize}
    \item The mass of each ring is \textit{small}, \(\Delta m\).
    \item The density of each ring is \textit{not small}, but depends on \(r\), so it's called \(\rho(r)\).
    \item The area of each ring is \textit{small} (thin!), but depends on \(r\), so let's call it \(\Delta A(r)\).
\end{itemize}
Therefore, \(\Delta m = \rho(r) \cdot \Delta A(r)\).

\item When we slice the satellite launch into small pieces along its journey: 
\begin{itemize}
    \item The energy needed to push the satellite through that slice is \textit{small}, \(\Delta E\).
    \item The force the satellite feels at that slice is \textit{not small} but depends on \(r\), \(F(r)\).
    \item The distance through that slice is \textit{small}, \(\Delta r\).
\end{itemize}
Therefore, \(\Delta E = F(r) \cdot \Delta r\).

\item When we slice the magnet push into small pieces along the wire:
\begin{itemize}
    \item The energy needed to push the magnet through this slice is \textit{small}: \(\Delta E\)
    \item The force the magnet feels at that slice is \textit{not small} but depends on \(x\), \(F(x)\).
    \item The distance through that slice is \textit{small}, \(\Delta x\).
\end{itemize}
Therefore, \(\Delta E = F(x) \cdot \Delta x\).

\item When we slice the atmosphere into small spherical shells:
\begin{itemize}
    \item The mass of the shell is \textit{small}, \(\Delta m\).
    \item The density in that shell is \textit{not small} but depends on \(h\), \(\rho(h)\).
    \item The volume of that shell is \textit{small} and depends on \(h\), \(\Delta V(h)\).
\end{itemize}
Therefore, \(\Delta m = \rho(h) \cdot \Delta h\).

\item When we slice the pool cue into particles: 
\begin{itemize}
    \item The mass of the rod-particle is \textit{small}, \(\Delta M\).
    \item The distance from the rod-particle to the ball is \textit{not small}, \(r\).
    \item The gravitational force the rod-particle exerts on the ball is \textit{small} (at least, relative to the force the whole rod ends up exerting), \(\Delta F\).
\end{itemize}
Therefore, \(\Delta F = \dfrac{G \cdot \Delta M \cdot m}{r^2}\). 

(I think I got the letters right for the mass of the rod vs. the mass of the ball.)

\end{enumerate}
%#####################################d############################
\end{document}